\chapter{Courbes multi-courbes et actualisation par le collatéral}\label{ch:rc:multi-curve-and-collateral-discounting}

Jusqu'à l'été 2007, un prêt interbancaire à trois mois et un prêt à trois mois
adossé à un collatéral coûtaient presque la même chose~: l'écart entre les deux était
d'environ dix points de base, et en euros moins de six. D'août 2007 à mai 2009, il est
resté au-dessus de cinquante points de base et, après la faillite de Lehman Brothers,
il a dépassé cent points de base pendant six mois. Tout outil de valorisation de swaps
utilisé dans la profession reposait sur une hypothèse que cette situation a rompue~:
une seule courbe projette les taux variables qu'un swap va payer et actualise chaque
flux de trésorerie, de sorte qu'une jambe variable indexée sur le taux à trois mois vaut
le pair. Dès que les taux interbancaires à trois et à six mois, ainsi que le taux au
jour le jour, portaient chacun leur propre prime de risque de crédit bancaire et de
liquidité, il n'existait plus de courbe unique. En quelques années, les dealers sont passés
à un cadre comportant une courbe par indice pour la projection et une courbe par convention
de collatéral pour l'actualisation, et en 2020 les chambres de compensation ont basculé
vers une nouvelle actualisation pour l'ensemble du marché en deux week-ends. Ce chapitre construit ce cadre à partir du constructeur de courbes du chapitre~1.

\section{La rupture de 2007 et la base de ténor}

Avant 2007, l'argument classique du livre~2, chapitre~9, tenait~: une jambe variable qui
verse le taux de chaque période, sur une courbe qui à la fois projette et actualise, vaut
$P(t_0)-P(t_n)$. Si une banque peut emprunter pendant six mois au taux Euribor à six mois
et le renouveler, une jambe variable à six mois vaut le pair, et les jambes à trois et à
six mois valent la même chose. Après 2007, un prêteur à six mois exigeait davantage que
deux prêts à trois mois consécutifs, en raison de la probabilité que l'emprunteur fasse
défaut au cours du second trimestre et de la valeur que représente la disponibilité
immédiate de ses liquidités. La différence est payée dans un swap de deux
jambes variables.

\begin{definition}[Base de ténor, swap de base de ténor]\label{def:rc:multi-curve-and-collateral-discounting:basis}
Un \emph{swap de base de ténor}\index{swap de base de ténor} (tenor basis swap) échange deux jambes variables
indexées sur des indices de ténors différents dans la même devise, par exemple l'Euribor à six mois
contre le taux au jour le jour capitalisé, avec un spread sur l'une des jambes qui rend le swap
de valeur nulle à l'origine. La \emph{base de ténor}\index{base de ténor} à une échéance
est ce spread au pair.
\end{definition}

\begin{proposition}[Base de ténor à partir de deux courbes de swaps]\label{prop:rc:multi-curve-and-collateral-discounting:basisrates}
Si un swap indexé sur le taux au jour le jour et un swap indexé sur le taux interbancaire, de même
échéance, ont des jambes fixes identiques et sont actualisés sur la même courbe, le spread de \omterm{def:rc:multi-curve-and-collateral-discounting:basis}{base
de ténor} (payé sur la jambe au jour le jour, selon l'échéancier de la jambe fixe) est égal à
la différence de leurs taux au pair, $b_n = S^{\text{IRS}}_n - S^{\text{OIS}}_n$.
\end{proposition}
\begin{proof}
Chaque swap au pair égalise sa jambe variable à $S_n A$, avec la même annuité $A$.
Recevoir la jambe interbancaire et payer la jambe au jour le jour plus $b$ vaut
$S^{\text{IRS}}_nA - S^{\text{OIS}}_nA - bA$, nul pour $b=S^{\text{IRS}}_n-S^{\text{OIS}}_n$.
\end{proof}

\section{Courbes d'actualisation et courbes de projection}

\begin{definition}[Courbe d'actualisation, courbe de projection, cadre multi-courbes]\label{def:rc:multi-curve-and-collateral-discounting:curves}
Une \emph{courbe d'actualisation}\index{courbe d'actualisation} donne la valeur actuelle d'un
flux de trésorerie payé selon une convention de collatéral donnée~; une \emph{courbe de
projection}\index{courbe de projection} d'un indice donne ses fixings à terme, sous la forme
$F(t;T_1,T_2) = \bigl(P^{\text{proj}}(t,T_1)/P^{\text{proj}}(t,T_2)-1\bigr)/\delta$,
sans servir à actualiser quoi que ce soit. Le \emph{cadre
multi-courbes}\index{cadre multi-courbes} valorise un swap en projetant chaque fixing variable
sur la courbe de projection de son indice, et en actualisant chaque flux de trésorerie sur la
courbe d'actualisation du collatéral de la transaction.
\end{definition}

Les «~facteurs d'actualisation~» de la \omterm{def:rc:multi-curve-and-collateral-discounting:curves}{courbe de projection} ne sont qu'un artifice pour stocker
des fixings à terme~; dans le langage du livre~4, chapitre~5, le terme forward de chaque
fixing est son espérance sous la probabilité forward de la courbe
d'\emph{actualisation}, ce qui explique pourquoi la \omterm{def:rc:multi-curve-and-collateral-discounting:curves}{courbe de projection} est calibrée avec
la \omterm{def:rc:multi-curve-and-collateral-discounting:curves}{courbe d'actualisation} maintenue fixe.

\begin{method}[Calibration en deux étapes]\label{met:rc:multi-curve-and-collateral-discounting:twostage}
(1) Calibrer la \omterm{def:rc:multi-curve-and-collateral-discounting:curves}{courbe d'actualisation} à partir des instruments indexés sur le taux au jour le
jour (voir la \cref{met:rc:curve-construction:newton}). (2) En la maintenant fixe, calibrer la \omterm{def:rc:multi-curve-and-collateral-discounting:curves}{courbe
de projection} de chaque indice à partir des fixings, des contrats à terme et des swaps de cet indice,
chacun valorisé comme $\sum_j F_j\delta_jP^{\text{disc}}(t_j)$ face à $K\sum_i
\delta_iP^{\text{disc}}(t_i)$. (3) Vérifier que chaque instrument se reprice et que la
base forward (fixing projeté moins le forward au jour le jour pour la
même période) est positive et régulière.
\end{method}

\omcode{code/firm/multicurve/firm_multicurve.py}{93}{115}{Calibration d'une courbe de projection sur une courbe d'actualisation fixe.}\label{lst:rc:multi-curve-and-collateral-discounting:proj}

\begin{example}[Un marché en euros]\label{ex:rc:multi-curve-and-collateral-discounting:euro}
Prenons des taux de swap €STR illustratifs allant de 1.95~\% à un an à 2.55~\% à dix ans
et 2.60~\% à trente ans, et des taux de swap Euribor à six mois supérieurs de 20 points de base à un an,
de 15 à dix ans et de 12 à trente ans, avec un fixing à six mois de 2.13~\%. La
\omterm{def:rc:multi-curve-and-collateral-discounting:curves}{courbe de projection} reprice les dix instruments Euribor~; par la
\cref{prop:rc:multi-curve-and-collateral-discounting:basisrates}, la \omterm{def:rc:multi-curve-and-collateral-discounting:basis}{base de ténor} vaut 20, 15 et 12
points de base à un, dix et trente ans. La base forward, fixing par fixing (la
\cref{fig:rc:multi-curve-and-collateral-discounting:forwards}), vaut 21.6 points de base sur la
première période de six mois et 11.8 sur la période qui commence dans dix ans~: une base
au pair est une moyenne des bases forward.
\end{example}

\begin{omfigure}
\begin{tikzpicture}
\begin{axis}[width=12cm,height=5.6cm,
  xlabel={début de la période de six mois (années)},ylabel={taux forward (\%)},
  tick label style={font=\small},label style={font=\small},
  xmin=0,xmax=30,ymin=1.8,ymax=3.4,grid=major,grid style={gray!25},
  legend columns=2,legend style={font=\footnotesize,at={(0.5,-0.25)},anchor=north,draw=none,fill=none,/tikz/every even column/.append style={column sep=8pt}},
  table/col sep=comma]
\addplot[omThm,very thick] table[x=t,y=euribor]{figdata/rates-credit-risk/02-multi-curve-and-collateral-discounting/forwards.csv};
\addplot[omDef,thick,dashed] table[x=t,y=ois]{figdata/rates-credit-risk/02-multi-curve-and-collateral-discounting/forwards.csv};
\legend{Forward Euribor à six mois (courbe de projection),Forward €STR à six mois (courbe d'actualisation)}
\end{axis}
\end{tikzpicture}

\omcaption{Taux forward pour des périodes de six mois issus des deux courbes en euros
de l'\cref{ex:rc:multi-curve-and-collateral-discounting:euro}. L'écart est la
base forward de ténor~: environ 22 points de base sur la première période, 12 après dix
ans, 9 à l'extrémité longue. Données~: les courbes illustratives du chapitre et le
tutoriel.}\label{fig:rc:multi-curve-and-collateral-discounting:forwards}
\end{omfigure}

\begin{example}[Un swap, deux cadres]\label{ex:rc:multi-curve-and-collateral-discounting:compare}
Un receveur de 3.50~\% contre l'Euribor à six mois pendant dix ans sur un notionnel de
100~millions d'euros, avec le swap Euribor à dix ans à 2.70~\%, vaut
7\,135\,408~euros dans le \omterm{def:rc:multi-curve-and-collateral-discounting:curves}{cadre multi-courbes} et 7\,075\,108~euros sur une
courbe unique construite à partir des swaps Euribor, qui actualise aux taux Euribor
plus élevés. Les deux cadres s'accordent sur le taux au pair, 2.70~\%, car chacun est
calibré sur lui~; ils diffèrent de 60\,300~euros sur un swap hors marché.
\end{example}

\section{Actualisation sur collatéral et credit support annex}

Pourquoi actualiser au taux au jour le jour~? Parce que c'est le taux versé sur les espèces
qui garantissent l'opération. Sous une credit support annex (livre~2,
chapitre~10), une partie dont le produit dérivé a une valeur positive $V_t$ reçoit
$V_t$ en espèces de la part de l'autre, et verse l'intérêt convenu sur ce montant~; le
produit est, en pratique, financé par son propre collatéral.

\begin{definition}[Taux de rémunération du collatéral, actualisation CSA]\label{def:rc:multi-curve-and-collateral-discounting:csa}
Le \emph{taux de rémunération du collatéral}\index{taux de rémunération du collatéral} $c_t$ d'une transaction est le taux
que verse le bénéficiaire du collatéral sur les espèces déposées, en vertu de la convention de collatéral de la transaction,
généralement l'indice de référence au jour le jour de la devise du collatéral.
L'\emph{actualisation CSA}\index{actualisation CSA} valorise chaque flux d'une transaction entièrement
garantie en espèces sur la courbe de son taux de rémunération du collatéral.
\end{definition}

\begin{proposition}[Le taux de rémunération du collatéral actualise]\label{prop:rc:multi-curve-and-collateral-discounting:csa}
Si une transaction est collatéralisée en continu et intégralement en espèces, rémunérées à $c_t$, et
que les couvertures qu'elle requiert sont financées au taux sans risque, sa valeur est
\[
V_t = \E^{\mathbb Q}_t\Bigl[\exp\Bigl(-\int_t^T c_u\,du\Bigr)V_T\Bigr].
\]
\end{proposition}
\begin{proof}[Preuve partielle]
Le détenteur de la transaction détient $V_t$ en collatéral, qu'il devra restituer, sur lequel il verse
$c_t$~; les positions de couverture rapportent et coûtent $r_t$ et s'annulent. Le compte
$V_t$ croît donc au taux $c_t$~: sous $\mathbb Q$, $dV_t = c_tV_t\,dt +
dM_t$ où $M$ est une martingale, dont la solution est la formule. Une preuve complète
(Piterbarg, 2010) détaille le financement des couvertures.
\end{proof}

\begin{omfigure}
\begin{tikzpicture}[font=\small,
  box/.style={draw,thick,rounded corners=2pt,align=center,minimum height=1.1cm,minimum width=3.2cm,inner sep=4pt},
  arr/.style={-{Stealth[length=2.5mm]},thick}]
\node[box,draw=omDef,fill=omDef!8] (bk) at (0,0) {banque\\(produit dérivé valant $V_t>0$)};
\node[box,draw=omThm,fill=omThm!8] (cp) at (7.6,0) {contrepartie\\(doit $V_t$)};
\draw[arr] (bk.15) -- node[above=6pt] {flux du swap} (cp.165);
\draw[arr] (cp.195) -- node[below=6pt] {flux du swap} (bk.345);
\draw[arr,omProp] (cp.north) to[bend right=25] node[above=3pt] {collatéral en espèces $V_t$ (marge de variation)} (bk.north);
\draw[arr,omExo] (bk.south) to[bend right=25] node[below=3pt] {intérêt $c_t V_t\,dt$ au taux de rémunération du collatéral} (cp.south);
\end{tikzpicture}

\omcaption{Un dérivé collatéralisé en espèces. La banque détient un collatéral égal
à la valeur de la transaction et verse sur celui-ci le \omterm{def:rc:multi-curve-and-collateral-discounting:csa}{taux de rémunération du collatéral}~; sa position dans la
transaction est financée par le collatéral, si bien que sa valeur croît au taux $c_t$ et est
actualisée sur la courbe de $c_t$
(la \cref{prop:rc:multi-curve-and-collateral-discounting:csa}).}\label{fig:rc:multi-curve-and-collateral-discounting:csa}
\end{omfigure}

Un swap compensé est identique, avec la chambre de compensation pour contrepartie~: la
marge de variation évolue chaque jour, et les intérêts qui s'y rattachent circulent dans l'autre sens.

\begin{definition}[Price alignment interest]\label{def:rc:multi-curve-and-collateral-discounting:pai}
Le \emph{price alignment interest}\index{price alignment interest} (PAI) est
l'intérêt que la chambre de compensation prélève auprès du bénéficiaire de la marge de variation en espèces,
et qu'elle verse au payeur, à un taux au jour le jour de son choix, de sorte qu'un swap compensé soit
économiquement équivalent à un swap bilatéral collatéralisé en espèces au même
taux. La chambre de compensation actualise au taux du PAI.
\end{definition}

\section{Choix du collatéral et collatéral multidevise}

De nombreuses conventions de collatéral permettent au déposant de choisir entre plusieurs devises ou
titres, et de les substituer ensuite.

\begin{definition}[Option de choix du collatéral]\label{def:rc:multi-curve-and-collateral-discounting:choice}
Une \emph{option de choix du collatéral}\index{option de choix du collatéral} est le
droit du déposant de choisir, puis de modifier, parmi plusieurs collatéraux éligibles, celui
qu'il remet. Exprimé dans la devise de la transaction, le déposant choisit
le collatéral qui rapporte le taux le plus élevé, de sorte que la transaction est actualisée à
$\max_i c^{(i)}_t$ sur les collatéraux éligibles.
\end{definition}

Le collatéral en devise étrangère rapporte le taux au jour le jour de cette devise~; converti
dans la devise de la transaction au moyen de la base de devises (livre~2,
chapitre~16), il rapporte le taux au jour le jour national plus ou moins la base.
Avec une base $b(t)$ pour prêter des euros contre des dollars, remettre des euros
rapporte $r_{\text{SOFR}}-b$ en dollars~: plus que des espèces en dollars lorsque $b<0$.

\begin{example}[Une transaction en dollars avec un CSA à deux devises]\label{ex:rc:multi-curve-and-collateral-discounting:choice}
Prenons une base illustrative de $-15$ points de base au court terme, croissant linéairement
jusqu'à $+5$ à dix ans. Les euros constituent le meilleur collatéral jusqu'à 7.5~ans, les dollars
ensuite. Un paiement de 100~millions de dollars dans dix ans vaut 68\,598\,292~dollars
dans un CSA libellé uniquement en dollars, et 68\,213\,510~dollars lorsque le payeur peut remettre l'une ou l'autre
devise~: le choix, utilisé à chaque date, relève le taux d'actualisation effectif
en moyenne de $56.25/10 = 5.6$ points de base par an et coûte au receveur
384\,782~dollars. Avec des taux et une base volatils, l'option vaut davantage que cette
valeur intrinsèque.
\end{example}

\begin{omfigure}
\begin{tikzpicture}
\begin{axis}[width=12cm,height=5.4cm,
  xlabel={temps (années)},ylabel={taux en dollars (\%)},
  tick label style={font=\small},label style={font=\small},
  xmin=0,xmax=10,ymin=3.1,ymax=4.4,grid=major,grid style={gray!25},
  legend columns=3,legend style={font=\footnotesize,at={(0.5,-0.25)},anchor=north,draw=none,fill=none,/tikz/every even column/.append style={column sep=8pt}},
  table/col sep=comma]
\addplot[omDef,thick] table[x=t,y=usd]{figdata/rates-credit-risk/02-multi-curve-and-collateral-discounting/collateral.csv};
\addplot[omProp,thick,dashed] table[x=t,y=eur]{figdata/rates-credit-risk/02-multi-curve-and-collateral-discounting/collateral.csv};
\addplot[omThm,line width=2.2pt,opacity=0.45] table[x=t,y=eff]{figdata/rates-credit-risk/02-multi-curve-and-collateral-discounting/collateral.csv};
\legend{espèces en dollars (SOFR),euros en espèces convertis en dollars,effectif (le maximum)}
\draw[gray,dashed] (axis cs:7.5,3.1) -- (axis cs:7.5,4.4);
\end{axis}
\end{tikzpicture}

\omcaption{\omterm{def:rc:curve-construction:pillar}{Taux forward instantanés} rémunérant deux collatéraux éligibles, en
dollars, et le taux d'actualisation effectif lorsque le déposant peut choisir. Les
courbes se croisent à 7.5 ans (pointillés), là où la base illustrative change de signe.
Données~: le tutoriel du chapitre.}\label{fig:rc:multi-curve-and-collateral-discounting:choice}
\end{omfigure}

\section{Après les taux interbancaires~: les changements d'actualisation de 2020}

Lorsque les indices de référence au jour le jour ont remplacé l'EONIA et le taux effectif des
federal funds en dollars comme taux de référence de leurs marchés, les chambres de
compensation ont modifié le taux de leur \omterm{def:rc:multi-curve-and-collateral-discounting:pai}{price alignment interest} (PAI), et avec lui la \omterm{def:rc:multi-curve-and-collateral-discounting:curves}{courbe
d'actualisation} de chaque swap compensé dans la devise.

\begin{definition}[Changement d'actualisation]\label{def:rc:multi-curve-and-collateral-discounting:switch}
Un \emph{changement d'actualisation}\index{changement d'actualisation} est une modification du
\omterm{def:rc:multi-curve-and-collateral-discounting:csa}{taux de rémunération du collatéral} d'un ensemble de transactions, et donc de leur \omterm{def:rc:multi-curve-and-collateral-discounting:curves}{courbe d'actualisation},
à une date fixée, les variations de valeur qui en résultent étant compensées en espèces et,
le cas échéant, les variations de risque étant compensées par des swaps de base.
\end{definition}

\begin{dated}{2026-09}{Les changements de 2020}\label{dat:rc:multi-curve-and-collateral-discounting:switches}
L'EONIA était calculé comme l'€STR plus 8.5 points de base du premier octobre 2019 jusqu'à sa
suppression, le 3 janvier 2022. LCH a fait passer l'actualisation et le price alignment
de tous les swaps en euros, de l'EONIA à l'€STR flat au 27 juillet 2020, en compensant chaque
compte en espèces pour la variation de la valeur actuelle nette, les flux projetés étant
maintenus fixes (une méthode «~constant-forward~»). Elle a fait passer les swaps en dollars du federal
funds vers le SOFR pendant le week-end du 16 au 19 octobre 2020, en espèces pour la variation
de valeur et des swaps de base federal funds/SOFR pour la variation de risque~: plus
d'un million de contrats et 120~billions de dollars de notionnel~; les clients qui ne voulaient pas
des swaps de base les ont fait mettre aux enchères, pour 24~milliards de dollars de notionnel net, auprès de 18 dealers
à un coût proche de zéro.
\end{dated}

Comme l'EONIA était fixé à l'€STR plus 8.5 points de base, le \omterm{def:rc:multi-curve-and-collateral-discounting:switch}{changement d'actualisation} en euros était
un décalage déterministe de la \omterm{def:rc:multi-curve-and-collateral-discounting:curves}{courbe d'actualisation}~: chaque facteur d'actualisation augmentait de
$e^{0.00085\,t}$. Un swap dans la monnaie, dont les flux futurs sont des encaissements nets,
gagnait de la valeur~; un swap hors de la monnaie en perdait~; la compensation
inversait chaque variation.

\omcode{code/rates-credit-risk/02-multi-curve-and-collateral-discounting/python/rc_multicurve.py}{125}{135}{Valorisation «~constant-forward~» du changement d'actualisation par la chambre de compensation.}\label{lst:rc:multi-curve-and-collateral-discounting:switch}

\begin{omfigure}
\begin{tikzpicture}
\begin{axis}[width=11.5cm,height=5.2cm,ybar,bar width=14pt,
  ylabel={variation de valeur (milliers d'EUR)},xlabel={échéance résiduelle},
  tick label style={font=\small},label style={font=\small},
  xmin=-0.6,xmax=7.6,ymin=0,ymax=650,grid=major,grid style={gray!25},
  xtick=data,xticklabels={2A,3A,5A,7A,10A,15A,20A,30A},
  nodes near coords,nodes near coords style={font=\scriptsize},
  table/col sep=comma]
\addplot[fill=omDef!55,draw=omDef] table[x=i,y=change]{figdata/rates-credit-risk/02-multi-curve-and-collateral-discounting/switch.csv};
\end{axis}
\end{tikzpicture}

\omcaption{Gain au \omterm{def:rc:multi-curve-and-collateral-discounting:switch}{changement d'actualisation} en euros, pour un receveur de 1.50~\% contre
l'Euribor à six mois sur 100~millions d'euros, par échéance résiduelle, sur des courbes
illustratives de juillet 2020 (taux des swaps en euros entre $-0.50~\%$ et $+0.08~\%$). La
compensation en espèces de la chambre de compensation a annulé chaque gain. Données~: le
tutoriel du chapitre.}\label{fig:rc:multi-curve-and-collateral-discounting:switch}
\end{omfigure}

\section{Tutoriel~: un marché en euros sur deux courbes}\label{tut:rc:multi-curve-and-collateral-discounting:tutorial}

\begin{tutorial}
\textbf{Objectif.} Construire une \omterm{def:rc:multi-curve-and-collateral-discounting:curves}{courbe d'actualisation} €STR et une \omterm{def:rc:multi-curve-and-collateral-discounting:curves}{courbe de projection} Euribor à six mois,
lire la \omterm{def:rc:multi-curve-and-collateral-discounting:basis}{base de ténor}, valoriser un swap dans les deux cadres, et valoriser le
\omterm{def:rc:multi-curve-and-collateral-discounting:switch}{changement d'actualisation} de 2020. \textbf{État final~:}
les \cref{fig:rc:multi-curve-and-collateral-discounting:forwards,fig:rc:multi-curve-and-collateral-discounting:switch}
et les chiffres de l'\cref{ex:rc:multi-curve-and-collateral-discounting:compare}.
\begin{tutsteps}
\item \textbf{\omterm{def:rc:multi-curve-and-collateral-discounting:curves}{Courbe d'actualisation}}~: \texttt{discount()}, neuf swaps €STR avec
le constructeur du chapitre~1, interpolation monotone convex.
\item \textbf{\omterm{def:rc:multi-curve-and-collateral-discounting:curves}{Courbe de projection}}~: \texttt{projection()} calibre le fixing et
neuf swaps Euribor avec la \omterm{def:rc:multi-curve-and-collateral-discounting:curves}{courbe d'actualisation} maintenue fixe (voir le \cref{lst:rc:multi-curve-and-collateral-discounting:proj})~; vérifier
\texttt{basis\_table()} renvoie 20, 19, 18, 17, 16, 15, 14, 13 et 12 points
de base.
\item \textbf{Deux cadres}~: \texttt{discounting\_comparison()}.
\item \textbf{Le \omterm{def:rc:multi-curve-and-collateral-discounting:switch}{changement d'actualisation}}~: \texttt{switch\_value(20, 0.015, 5e8)}~;
\texttt{fig\_rc\_multicurve.py} écrit les graphiques.
\end{tutsteps}
\textbf{Pour aller plus loin.} Rendre la base plate à 15 points de base et vérifier
que la base forward devient plate elle aussi~; valoriser le \omterm{def:rc:multi-curve-and-collateral-discounting:switch}{changement d'actualisation} avec la \omterm{def:rc:multi-curve-and-collateral-discounting:curves}{courbe de projection}
recalibrée sur la \omterm{def:rc:multi-curve-and-collateral-discounting:curves}{courbe d'actualisation} €STR au lieu d'être maintenue fixe, et
comparer avec la méthode de la chambre de compensation.
\end{tutorial}

\section{Construction~: courbes d'actualisation et de projection}\label{bld:rc:multi-curve-and-collateral-discounting:multicurve}

\begin{build}
\bfield{Rôle} Chaque instrument de taux du livre est valorisé avec une \omterm{def:rc:multi-curve-and-collateral-discounting:curves}{courbe de projection}
pour son indice et une \omterm{def:rc:multi-curve-and-collateral-discounting:curves}{courbe d'actualisation} pour son collatéral.
\bfield{Interface} \texttt{IborSwap}, \texttt{Fixing}~;
\texttt{calibrate\_projection(spot, disc, instruments, kind)}~;
\texttt{ibor\_swap\_pv(proj, disc, \ldots)}~; \texttt{tenor\_basis}~;
\texttt{ShiftedCurve(base, shift)}~; \texttt{CollateralChoiceCurve(base,
spreads)}~; toutes les courbes exposent \texttt{df(date)}.
\bfield{Règles} La \omterm{def:rc:multi-curve-and-collateral-discounting:curves}{courbe d'actualisation} est calibrée en premier et jamais modifiée par la
calibration de projection~; jambes fixes annuelles ACT/360, jambes variables sur le
ténor de l'indice~; choix du collatéral comme maximum ponctuel des taux éligibles~;
\texttt{firm\_curvebuild} et le \texttt{firm\_curve} du livre~2 importés, non modifiés.
\bfield{Tests d'acceptation} \texttt{code/firm/multicurve/tests/}~: la \omterm{def:rc:multi-curve-and-collateral-discounting:curves}{courbe de projection}
revalorise ses swaps sur sa \omterm{def:rc:multi-curve-and-collateral-discounting:curves}{courbe d'actualisation}~; la \omterm{def:rc:multi-curve-and-collateral-discounting:basis}{base de ténor} est égale à la
différence des taux au pair~; sans base, les forwards de projection et d'actualisation coïncident~;
les valeurs payeur et receveur sont opposées~; un choix du collatéral n'augmente jamais la
valeur d'une créance et se réduit à la courbe unique avec un seul collatéral.
\bfield{Pour aller plus loin} Plusieurs indices (un, trois, six, douze mois) à partir de swaps
de base~; une \omterm{def:rc:multi-curve-and-collateral-discounting:curves}{courbe d'actualisation} multidevise à partir de swaps de base de devises~; le
choix du collatéral avec base volatile par simulation.
\end{build}

\begin{omsources}
\item V.\ Piterbarg, ``Funding beyond discounting: collateral agreements and
derivatives pricing'', \emph{Risk}, février 2010~; ``Cooking with collateral'',
\emph{Risk}, 2012.
\item A.\ De Socio, ``The interbank market after the financial turmoil'', Banca
d'Italia, Temi di Discussione 819, 2011.
\item LCH, \emph{SwapClear transition to €STR discounting}; \emph{SOFR
discounting: plan for the SwapClear compensation process}, 2019--2020.
\item M.\ Henrard, \emph{Interest Rate Modelling in the Multi-Curve Framework},
Palgrave Macmillan, 2014.
\end{omsources}

\section{Exercices}

\begin{exercise}[$\star$]\label{exo:rc:multi-curve-and-collateral-discounting:1}
Le swap à dix ans contre l'Euribor six mois est à 2.70~\% et le swap à dix ans
€STR est à 2.55~\%, à jambes fixes annuelles identiques. Quelle est la \omterm{def:rc:multi-curve-and-collateral-discounting:basis}{base de ténor} à dix ans
et sur quelle jambe est-elle versée~?
\end{exercise}

\begin{exercise}[$\star$]\label{exo:rc:multi-curve-and-collateral-discounting:2}
Avant 2007, le spread à trois mois entre taux non garanti et taux garanti était
d'environ dix points de base~; ensuite, il a dépassé cinquante. Expliquer en deux phrases
pourquoi cela a mis fin au cadre à courbe unique.
\end{exercise}

\begin{exercise}[$\star$]\label{exo:rc:multi-curve-and-collateral-discounting:3}
Un swap en euros est collatéralisé en espèces en euros rémunérées au €STR. Quelle courbe l'actualise,
quelle courbe projette ses fixings à six mois, et que change-t-il si le CSA rémunère
le €STR majoré de 10 points de base~?
\end{exercise}

\begin{exercise}[$\star\star$]\label{exo:rc:multi-curve-and-collateral-discounting:4}
Dans l'\cref{ex:rc:multi-curve-and-collateral-discounting:choice}, calculer l'intégrale
du spread effectif sur dix ans et le rapport des facteurs d'actualisation qui en
résulte, et vérifier le résultat de 384\,782~USD.
\end{exercise}

\begin{exercise}[$\star\star$]\label{exo:rc:multi-curve-and-collateral-discounting:5}
Expliquer les 60\,300~EUR de
l'\cref{ex:rc:multi-curve-and-collateral-discounting:compare}~: dans quel sens se trompe la courbe unique,
et pourquoi les deux cadres sont-ils justes sur le taux au pair~?
\end{exercise}

\begin{exercise}[$\star\star$]\label{exo:rc:multi-curve-and-collateral-discounting:6}
Au basculement vers l'euro, un client \emph{verse} 1.50~\% sur 100~millions d'euros pendant
dix ans. En utilisant la \cref{fig:rc:multi-curve-and-collateral-discounting:switch},
donner la variation de sa valeur et préciser qui verse la compensation à qui.
\end{exercise}

\begin{exercise}[$\star\star\star$]\label{exo:rc:multi-curve-and-collateral-discounting:7}
\emph{Programmation.} Avec \texttt{forward\_table()}, donner la base forward sur la
période de six mois qui débute dans dix ans, la comparer à la base au pair à dix ans de
15 points de base, et expliquer l'écart.
\end{exercise}

\begin{exercise}[$\star\star\star$]\label{exo:rc:multi-curve-and-collateral-discounting:8}
\emph{Trouver la faille.} «~Notre CSA permet à la contrepartie de déposer des euros ou des dollars~;
aujourd'hui les dollars sont moins chers à déposer, donc l'\omterm{def:rc:multi-curve-and-collateral-discounting:choice}{option de choix du collatéral} est hors de la monnaie et nous
pouvons valoriser la transaction sur la courbe en dollars~». Corriger-la.
\end{exercise}

\section{Problème~: le week-end du basculement}

\begin{problem}[{Problème du week-end --- le swap d'un fonds de pension au changement d'actualisation en euros}]\label{pb:rc:multi-curve-and-collateral-discounting:1}
Le vendredi 24 juillet 2020, un fonds de pension reçoit annuellement 1.50~\% contre
l'Euribor six mois sur 500~millions d'euros, compensé auprès de LCH, avec vingt ans restants~;
il a été conclu en 2010. Sur les courbes illustratives de juillet 2020 du chapitre
(swaps €STR de $-0.50~\%$ à $-0.05~\%$, swaps Euribor de 12 à 17 points de base
plus élevés), le taux au pair Euribor à vingt ans est de 0.08~\%. Le lundi 27 juillet, la
chambre de compensation actualise au €STR flat au lieu de l'EONIA ($=$ €STR $+$ 8.5
points de base).

\medskip\noindent\textbf{Partie I --- La position.}
\begin{enumerate}
\item Le swap est-il un actif ou un passif du fonds, et pourquoi~?
\item Donner sa valeur le vendredi sous l'actualisation à l'EONIA.
\item Quels flux du swap sont connus, lesquels sont projetés~?
\item Que reçoit le fonds au titre de sa marge de variation, avant et après le basculement~?
\item Pourquoi le taux d'actualisation d'un swap compensé est-il égal au taux PAI~?
\end{enumerate}

\medskip\noindent\textbf{Partie II --- Le basculement.}
\begin{enumerate}[resume]
\item De combien chaque facteur d'actualisation jusqu'à $t$ années change-t-il~?
\item Donner la valeur du swap le lundi selon la méthode du forward constant de la chambre de compensation.
\item Donner la variation de valeur. Qui y gagne~?
\item Qui verse la compensation en espèces, pour quel montant, et à qui~?
\item Pourquoi maintenir fixes les flux projetés plutôt que recalibrer la
\omterm{def:rc:multi-curve-and-collateral-discounting:curves}{courbe de projection} sur la nouvelle \omterm{def:rc:multi-curve-and-collateral-discounting:curves}{courbe d'actualisation}~?
\end{enumerate}

\medskip\noindent\textbf{Partie III --- Le risque.}
\begin{enumerate}[resume]
\item Après le basculement, combien le swap perd-il si la \omterm{def:rc:multi-curve-and-collateral-discounting:curves}{courbe d'actualisation} €STR
monte d'un point de base, les projections étant maintenues~?
\item Avant le basculement, ce même chiffre décrivait le risque EONIA. Quel risque a
changé, et pourquoi le basculement en euros n'a-t-il pas exigé de compensation de swap de base alors que le
basculement en dollars en a exigé une~?
\item Un payeur de 1.50~\% aurait-il gagné ou perdu au basculement~?
\item Pourquoi le gain du receveur à trente ans est-il plus du double de celui du receveur à quinze ans
(la \cref{fig:rc:multi-curve-and-collateral-discounting:switch})~?
\item Le basculement modifie-t-il la position économique du fonds~? Ses comptes~?
\end{enumerate}

\medskip\noindent\textbf{Partie IV --- Jugement.}
\begin{enumerate}[resume]
\item Pourquoi le marché a-t-il fait basculer tous les swaps compensés en un seul week-end plutôt que
chaque transaction à son propre rythme~?
\item Qu'aurait exigé une transaction bilatérale sous CSA rémunéré à l'EONIA~?
\item Un dealer aurait-il pu tirer profit de la connaissance de la date du basculement~? Pourquoi pas beaucoup~?
\item Énoncer le \emph{résultat nommé}~: la variation de valeur du fonds et la compensation,
et sa sensibilité à l'actualisation après le basculement.
\item En une phrase~: qu'est-ce qui détermine la \omterm{def:rc:multi-curve-and-collateral-discounting:curves}{courbe d'actualisation} d'un produit dérivé~?
\end{enumerate}
\end{problem}

\section{Questions d'entretien}

\begin{interviewq}[$\star$ \iqroles{chercheur, banque}]\label{iq:rc:multi-curve-and-collateral-discounting:1}
Pourquoi actualise-t-on les swaps collatéralisés au taux au jour le jour et non
à l'indice qu'ils versent~?
\end{interviewq}

\begin{interviewq}[$\star\star$ \iqroles{chercheur, développeur}]\label{iq:rc:multi-curve-and-collateral-discounting:2}
Expliquez-moi comment construire une courbe Euribor à six mois après 2008.
Qu'est-ce qui est calibré en premier, et pourquoi~?
\end{interviewq}

\begin{interviewq}[$\star\star$ \iqroles{trader}]\label{iq:rc:multi-curve-and-collateral-discounting:3}
Qu'est-ce qu'un \omterm{def:rc:multi-curve-and-collateral-discounting:basis}{swap de base de ténor}, et pourquoi la base à six mois est-elle positive~?
\end{interviewq}

\begin{interviewq}[$\star\star$ \iqroles{banque, risque}]\label{iq:rc:multi-curve-and-collateral-discounting:4}
Un CSA autorise le collatéral dans trois devises. Comment cela modifie-t-il la valeur
des transactions qui y sont soumises, et dans quel sens~?
\end{interviewq}

\begin{interviewq}[$\star\star\star$ \iqroles{chercheur, risque}]\label{iq:rc:multi-curve-and-collateral-discounting:5}
Une chambre de compensation modifie son taux d'actualisation. Que devient la valeur
et le risque d'un portefeuille de swaps, et comment compenseriez-vous équitablement les membres~?
\end{interviewq}

\begin{interviewq}[$\star\star\star$ \iqroles{développeur}]\label{iq:rc:multi-curve-and-collateral-discounting:6}
Votre pricer actualise par erreur sur la \omterm{def:rc:multi-curve-and-collateral-discounting:curves}{courbe de projection}. Quelles transactions
révèlent l'erreur, et quelle en est l'ampleur~?
\end{interviewq}
